\section{INTRODUCTION}
\label{sec:introduction}
The practical appeal of O-ISAC lies in reuse. Optical networks already deploy
sources, propagation paths, receivers, and signal processors, and these
resources can reveal information about the environment while carrying data.
Distributed sensing over fiber, positioning based on visible light, and
turbulence monitoring illustrate the breadth of this opportunity
\cite{OISAC_SCR00013,OISAC_SCR00196,OISAC_SCR01085}.

Sharing turns that opportunity into a systems question. A resource or receiver
chosen for one function can change the information available to the other.
The central engineering problem is therefore not whether light can support
communication and sensing, but what can be shared and how that sharing changes
system design.

This question has gained importance since integrated sensing and communication
(ISAC) was identified as one of the usage scenarios for International Mobile
Telecommunications for 2030 and beyond (IMT-2030)
\cite{ITUR_M2160_2023}. Optical implementations span guided fiber, free space
optical and visible light communication links, systems in the millimeter wave
and terahertz bands enabled by photonics, and hybrid designs
\cite{OISAC_SCR00916,OISAC_SCR00036,OISAC_SCR00083}. Their propagation paths,
observables, hardware, and performance measures differ. O-ISAC is therefore a
family of related architectures rather than one interchangeable platform.

That diversity requires both a clear survey boundary and care in comparison.
Here, O-ISAC covers optical or photonic systems in which communication and
sensing are considered, integrated, designed, optimized, or evaluated together.
Their common technical setting may be an architecture, optical link, signal or
resource framework, hardware platform, channel model, or application scenario.
Sharing a waveform or operating concurrently are possible forms of coupling,
not inclusion requirements. Within this scope, familiar metric names are most
informative when read together with their measurement and operating context.
Preserving that context allows results from different platforms to be
interpreted consistently and compared where their defining conditions align.

\subsection{Prior Syntheses and Survey Positioning}
\label{sec:related_surveys}

Existing surveys illuminate complementary parts of O-ISAC, but differ in the
evidence they organize and the questions they help readers answer.
Table~\ref{tab:prior_surveys} organizes them by reader task, evidence unit, and
comparison boundary.

% Integration requirement: the manuscript driver must load array, booktabs,
% and tabularx before this section is included.
\begin{table*}[!t]
\caption{Prior O-ISAC syntheses by reader task, evidence unit, and comparison
boundary. Families support navigation rather than quality ranking.}
\label{tab:prior_surveys}
\centering
\begingroup
\footnotesize
\setlength{\tabcolsep}{2.2pt}
\renewcommand{\arraystretch}{1.00}
\begin{tabularx}{\textwidth}{@{}%
>{\raggedright\arraybackslash}p{0.112\textwidth}
>{\raggedright\arraybackslash}p{0.140\textwidth}
>{\raggedright\arraybackslash}p{0.120\textwidth}
>{\raggedright\arraybackslash}p{0.190\textwidth}
>{\raggedright\arraybackslash}p{0.160\textwidth}
>{\raggedright\arraybackslash}X@{}}
\toprule
Synthesis family & Sources & Optical scope & Evidence unit and comparison logic &
Reader task & Boundary \\
\midrule
Architecture and system design &
\cite{Wen2024OISACArchitectures,Su2025OISACTheoryPractice,
Chu2025ROISACOverview,Wang2026CrossLayerISACoF} &
General O-ISAC, free space optical links, retroreflective links, and fiber
networks &
Architecture, optimization, target, or link model; taxonomy and relations
within each scenario &
How are communication and sensing coupled? &
Relations remain architecture and assumption specific. \\
\addlinespace[1pt]
Analytical limits and estimation methods &
\cite{Khorasgani2026GeneralOpticalISACSurvey,Rode2025PolarizationISAC} &
General optical systems and polarization sensing in coherent fiber links &
Model, bound, objective, estimator, or learning method; comparison within a
model and among methods &
What limits follow from the assumptions? &
A single model does not define an empirical scale for an entire modality. \\
\addlinespace[1pt]
Technological evolution, standardization, and deployment &
\cite{Zhang2026ISACEvolution,Liu2025OpticalNetworkISACProspect,
Cao2026SubmarineFiberISAC,Zhang2026OpticalTransmissionNetworks,
Ip2026DeployedTelecomCableISAC} &
Evolution from radio to optical systems; terrestrial and submarine networks &
Historical phase, standardization path, network architecture, or deployment;
synthesis by period and ecosystem &
How has O-ISAC moved toward deployment? &
A shared trajectory does not make infrastructures commensurate. \\
\addlinespace[1pt]
Physical media and optical networks &
\cite{Liang2024LiSACReview,Zhang2023OpticalNetworkISACFrontier,
He2025FiberISACAdvances,Liu2026OpticalFiberISACReview,
Lu2026ConvergedFiberISACReview,Wang2026FiberNetworkISAC} &
Visible light ISAC, guided fiber systems, and optical networks &
Medium, observable, sensing path, or network configuration; taxonomy within
related solution classes &
How does each path support sensing? &
Transfer requires aligned media, observables, paths, and conditions. \\
\addlinespace[1pt]
Waveforms, resource multiplexing, and photonic hardware &
\cite{Yang2025DSCMISACReview,He2026PhotonicFiberISAC,
Yu2023PhotonAssistedTHz,Bai2025MicrowavePhotonicsISAC,
Wang2025MicrowavePhotonicsISAC,Lyu2026PhotonicTHzWaveforms} &
Fiber, microwave photonics, and millimeter wave or terahertz systems enabled by
photonics &
Waveform, multiplexing, front end, or processing chain; mechanism taxonomy with
selected demonstrations &
Which mechanisms enable joint operation? &
Shared photonics does not equate paths, observables, or metrics. \\
\addlinespace[1pt]
Applications and challenges across platforms &
\cite{Mohsan2026OISACReview} &
Visible light communication, free space optical, fiber, photonics, and hybrid
settings &
Application, use case, solution class, or challenge; thematic synthesis across
platforms &
Which applications and gaps recur? &
Breadth connects research lines without making results commensurate. \\
\bottomrule
\end{tabularx}
\endgroup
\end{table*}

Across these perspectives, the remaining task is to relate primary evidence
without separating a result from the setting in which it acquired meaning.
This article is a survey of 206 unique O-ISAC studies in the primary evidence
window from 1 January 2020 to 22 June 2026. Its central objective is to organize
and compare that evidence while preserving physical, measurement, and validation
context.

The survey makes four contributions:

\begin{itemize}
\item A taxonomy separates physical modality from the mechanisms that couple
communication and sensing.

\item A context-preserving comparison framework keeps each performance result
connected to its measurement definition and operating conditions.

\item A condition-aware synthesis brings recurring communication and sensing
tradeoff mechanisms together while retaining the conditions that shape their
interpretation.

\item An evidence-maturity assessment links validation, benchmark readiness,
and reproducibility to research priorities grounded in the reviewed evidence.
\end{itemize}

The organization of the article follows the analytical sequence of the survey.
Sections~\ref{sec:foundations_comparison} and~\ref{sec:review_methods}
establish the comparison framework across optical platforms and describe the
review methodology and evidence base. Sections~\ref{sec:optical_modality_families}
through~\ref{sec:technologies_applications_6g} present the technical synthesis
of optical platforms, integration architectures, performance metrics, joint
design tradeoffs, validation, enabling technologies, and applications.
Section~\ref{sec:discussion_roadmap} discusses the research priorities and
limitations that emerge from this synthesis, and
Section~\ref{sec:conclusion} concludes the article.
